\vfill\pagebreak
\section{Two names for the same elements}
\label{sec:mutability:sharing}

Assigning a variable to another copies its value. What that copies depends on
where the value is. The next listing assigns an array, then a slice, and
changes the first variable afterwards.

\lstinputlisting[style=coloredverbatim, caption={Assigning an array and a slice, \textit{examples/mutability/share.yr}}, label=lst:mutability:share]{examples/mutability/share.yr}
\vspace{-10pt}%
\begin{lstlisting}[style=bashVerb]
[31, 29, 31] [31, 28, 31]
[7, 3, 5] [7, 3, 5]
\end{lstlisting}

The array behaves as an integer would. \token{let before = days;} on line~5
copies the three elements into \token{before}, and line~6 changes
\token{days[1]} only: \token{before} keeps the 28 it received.

The slice does not. \token{let seen = primes;} on line~10 copies what
\token{primes} holds, which is a length and an address. \token{seen} designates
the same elements as \token{primes}, and when line~11 changes the first of them,
\token{seen[0]} is 7 as well. Figure~\ref{fig:mutability:shared} shows both
cases.

\begin{figure}[h]
  \centering
  \begin{adjustbox}{max width=\linewidth}
    \begin{tikzpicture}
      \foreach \value [count=\i from 0] in {31, 29, 31} {
        \ifnum\i=1
          \node[memhit] (d\i) at (\i * 1.0, 0) {\value};
        \else
          \node[memcell] (d\i) at (\i * 1.0, 0) {\value};
        \fi
      }
      \node[cflab, left=3mm of d0] (dn) {\token{days}};
      \foreach \value [count=\i from 0] in {31, 28, 31} {
        \node[memcell] (b\i) at (\i * 1.0, -1.0) {\value};
      }
      \node[cflab, left=3mm of b0] {\token{before}};

      \node[memcell] (plen) at (0, -2.2) {3};
      \node[memaddr, right=0mm of plen] (pptr) {};
      \node[cflab, left=3mm of plen] (pn) {\token{primes}};
      \node[memcell] (slen) at (0, -3.2) {3};
      \node[memaddr, right=0mm of slen] (sptr) {};
      \node[cflab, left=3mm of slen] (sn) {\token{seen}};
      \node[memframe, fit=(dn)(d2)(sn)(sptr)] (main) {};
      \node[memtitle] at ([xshift=2mm]main.north west) {main};

      \foreach \value [count=\i from 0] in {7, 3, 5} {
        \ifnum\i=0
          \node[memhit] (h\i) at (5.2 + \i * 1.0, -2.7) {\value};
        \else
          \node[memcell] (h\i) at (5.2 + \i * 1.0, -2.7) {\value};
        \fi
      }
      \draw[mempoint] (pptr.center) to[out=0, in=180] (h0.west);
      \draw[mempoint] (sptr.center) to[out=0, in=180] (h0.west);

      \node[memregion] (stack) at ($(main.north) + (0, 5mm)$) {stack};
      \node[memregion] at (h1 |- stack.base) {heap};
      \draw[memsep] (4.3, 0.9) -- (4.3, -4.0);
    \end{tikzpicture}
  \end{adjustbox}
  \caption{\label{fig:mutability:shared} The variables of
    Listing~\ref{lst:mutability:share} at its end. \token{before} is a copy of
    the elements of \token{days}, and did not follow the change of
    \token{days[1]}. \token{seen} is a copy of the length and the address of
    \token{primes}: both designate the same elements, and \token{primes[0] = 7;}
    changed what both hold.}
\end{figure}


Sharing is what makes slices cheap to pass around: assigning one or passing it
to a function copies two numbers, whatever its length. The functions
\token{sum} and \token{average} of Listing~\ref{lst:collections:average}
received the slice of their caller that way, and did not copy its elements. A
part of a slice, \token{s[1 .. 4]}, shares the elements of \token{s} in the same
way (Section~\ref{sec:collections:slices}).

As long as no one changes the elements, sharing them cannot be seen: two
variables that designate the same elements \token{[2, 3, 5]} and two variables
that hold a copy each print the same thing. It shows once one of them changes
an element, as \token{primes} does on line~11. \token{seen} is a constant, and
cannot change the elements itself, but it sees the changes made through
\token{primes}. The next sections explain which variables may change shared
elements, and how a program says that it expects them to.

\subsection{Giving a variable a new slice}

Changing an element and giving the variable a new value are two different
things. Changing an element writes into the block that the slice designates, and
every variable that designates it sees the change. Giving the variable a new
value replaces the length and the address that it holds, and the others are not
affected.

\begin{lstlisting}[style=coloredverbatim]
let mut primes = copy [2, 3, 5];
let seen = primes;
primes = copy [11, 13];
println(primes, " ", seen);
\end{lstlisting}
\vspace{-10pt}%
\begin{lstlisting}[style=bashVerb]
[11, 13] [2, 3, 5]
\end{lstlisting}

After the third line, \token{primes} designates a new block of two elements,
and \token{seen} still designates the first one.

\tokennolst{\mtildeascii{}=} (Section~\ref{sec:collections:slices}) also gives
the variable a new slice, but not always a new block. When the block of the
slice has room after its last element, the elements appended are written there,
and the variable designates the same block, one element longer. A variable that
shared the elements before still shares the first ones.

\begin{lstlisting}[style=coloredverbatim]
let dmut numbers = copy [1, 2, 3];
let seen = numbers;
numbers ~= [4];
numbers[0] = 9;
println(numbers, " ", seen);
\end{lstlisting}
\vspace{-10pt}%
\begin{lstlisting}[style=bashVerb]
[9, 2, 3, 4] [9, 2, 3]
\end{lstlisting}

Whether the block has room depends on how it was allocated, and a program
should not count on either outcome. When two variables must not share anything,
one of them is made with \token{copy}, which Section~\ref{sec:mutability:copies}
presents.
